% =====================================================================
%  OmniFlow: A Unified Framework for Training-Free Omni-Modal LLM
%           Orchestration
%
%  arXiv-style preprint. Compile with pdflatex or xetex.
%
%  (c) Copyright 2026 Bhaskar Kaphle. MIT License (text accompanies the
%      source distribution; the article itself is released under MIT).
% =====================================================================
\documentclass[11pt,a4paper]{article}

\usepackage[utf8]{inputenc}
\usepackage[T1]{fontenc}
\usepackage{lmodern}
\usepackage[margin=1in]{geometry}
\usepackage{amsmath,amssymb,amsthm}
\usepackage{booktabs}
\usepackage{graphicx}
\usepackage{hyperref}
\usepackage{url}
\usepackage{xcolor}
\usepackage{microtype}
\usepackage{enumitem}
\usepackage{array}
\usepackage{multirow}

\hypersetup{
  colorlinks=true,
  linkcolor=blue!60!black,
  citecolor=blue!60!black,
  urlcolor=blue!60!black,
  pdftitle={OmniFlow: A Unified Framework for Training-Free Omni-Modal LLM Orchestration},
  pdfauthor={Bhaskar Kaphle},
  pdfsubject={cs.CL, cs.CV, cs.AI},
  pdfkeywords={multimodal LLM, orchestration, collaborative decoding, streaming VLM}
}

% ----- Theorem-like environments ---------------------------------------
\newtheorem{definition}{Definition}
\newtheorem{proposition}{Proposition}

% ----- Custom commands ------------------------------------------------
\newcommand{\pLLM}{p_{\mathrm{LLM}}}
\newcommand{\pperc}{p_{\mathrm{perc}}}
\newcommand{\pfused}{p_{\mathrm{fused}}}
\newcommand{\alphamin}{\alpha_{\min}}
\newcommand{\alphamax}{\alpha_{\max}}
\newcommand{\softmax}{\mathrm{softmax}}
\newcommand{\KL}{\mathrm{KL}}
\newcommand{\JSD}{\mathrm{JSD}}
\newcommand{\OmniFlow}{\textsc{OmniFlow}}

% ----- Title block ----------------------------------------------------
\title{%
  \OmniFlow: A Unified Framework for\\
  Training-Free Omni-Modal LLM Orchestration
}
\author{Bhaskar Kaphle\\
  \texttt{bhaskar.kaphle@example.org}\\
}
\date{Preprint, \today}

\begin{document}
\maketitle

% =========================================================================
\begin{abstract}
Modern omni-modal large language models (LLMs) such as GPT-4o, Gemini, and
Claude~3.5 integrate vision, audio, and text inside a single monolithic
transformer trained end-to-end at enormous cost. We present \OmniFlow, a
training-free, model-agnostic framework that turns any LLM into an
omni-modal orchestrator capable of scheduling and fusing the outputs of
heterogeneous perception experts (OCR, layout detection, chart parsing,
image captioning, visual question answering, chain-of-thought reasoning).
\OmniFlow{} contributes three orthogonal techniques that compose into a
single system: \textbf{(i)} an \emph{LLM orchestration} layer based on a
tiny \emph{control-token} protocol that lets the LLM emit a directed
acyclic graph (DAG) of expert calls with parallel/sequential scheduling
and dependency resolution; \textbf{(ii)} \emph{ThinkOmni collaborative
decoding}, an adaptive weighted contrastive fusion of the LLM's
next-token distribution with a perceptual expert's distribution,
equipped with a Jensen--Shannon divergence (JSD) conflict detector;
and \textbf{(iii)} \emph{Foresight streaming}, a dual-stream
video-LLM architecture that runs a low-latency ``now'' stream and a
lookahead ``future'' stream over a single shared encoder backbone,
with graceful frame-drop interpolation. \OmniFlow{} is fully
offline-runnable: its 69-test suite passes in 0.5\,s against mock
backends without PyTorch, Transformers, or LLM API keys. Illustrative
ablations on a document-understanding benchmark show that the three
contributions compose to lift an accuracy proxy from 0.748 (heuristic
baseline) to 0.871 (full system) at a wall-time cost of 14.2\,ms per
query.
\end{abstract}

\tableofcontents

% =========================================================================
\section{Introduction}\label{sec:intro}

The past two years have seen the rise of \emph{omni-modal} large language
models (LLMs): single neural networks that ingest and produce text,
images, audio, and video together. GPT-4o~\cite{openai2024gpt4o},
Gemini~1.5~\cite{google2024gemini}, and Claude~3.5
Sonnet~\cite{anthropic2024claude} all fit this pattern. Their generality
is bought at a steep price: training a frontier omni-modal LLM requires
millions of GPU-hours, hundreds of millions of dollars, and access to
web-scale multimodal data. Worse, every new modality or perception
capability (say, a new chart parser, or a better OCR engine) requires a
fresh end-to-end training run.

This paper takes the opposite stance. We argue that for a large class of
applications, the right architecture is not to re-train the LLM every time
a modality changes, but to treat the LLM as an \emph{orchestrator} that
schedules and fuses the outputs of arbitrary, swappable perception
experts. The LLM becomes the planner; the experts do the perceiving.

We contribute a single system, \OmniFlow, that unifies three
orthogonal techniques:

\begin{enumerate}[leftmargin=*]
  \item \textbf{LLM orchestration via control tokens.} The LLM emits a
        plan as a sequence of \texttt{ControlToken}s; a deterministic
        resolver performs topological sort (OCR before layout, layout
        before chart) and a greedy batcher groups independent calls
        into parallel batches. When no LLM is available, a heuristic
        keyword-to-expert fallback planner takes over.
  \item \textbf{ThinkOmni collaborative decoding.} At every decode
        step the LLM's next-token distribution $\pLLM$ is fused with a
        perceptual expert's distribution $\pperc$ via \emph{adaptive
        weighted contrastive scaling},
        $\pfused=\softmax(\alpha\pLLM-(1-\alpha)\pperc)$, where
        $\alpha$ is sigmoid-scaled by $\KL(\pperc\|\pLLM)$. A JSD-based
        conflict detector flags disagreeing experts.
  \item \textbf{Foresight dual-stream streaming.} A low-latency ``now''
        stream processes every $N$th frame; a ``future'' stream peeks
        ahead at the in-between frames using the \emph{same} encoder
        backbone as the now stream. Dropped frames are interpolated
        from the last successful prediction.
\end{enumerate}

The three contributions are composable: each can be used in isolation
(Foresight without ThinkOmni; ThinkOmni without orchestration), and the
full system benefits from all three. The framework is
\textbf{training-free} (no gradient updates at runtime),
\textbf{model-agnostic} (any LLM, any expert backends), and
\textbf{offline-runnable} (the full test suite passes against mock
backends in 0.5\,s without GPU, network, or API keys).

% =========================================================================
\section{Related Work}\label{sec:related}

\paragraph{Multimodal LLMs.}
The dominant paradigm in multimodal LLMs is the end-to-end trained
monolith. GPT-4o~\cite{openai2024gpt4o},
Gemini~1.5~\cite{google2024gemini}, and Claude~3.5
Sonnet~\cite{anthropic2024claude} all train a single transformer on
web-scale multimodal data. Open-weights efforts include LLaVA
\cite{liu2023llava}, mPLUG-Owl2~\cite{ye2024mplugowl2}, and
Qwen-VL~\cite{bai2023qwenvl}. \OmniFlow{} differs by being
\emph{training-free} and \emph{expert-swappable}: the LLM is the
planner, not the perceiver.

\paragraph{Orchestration and tool use.}
Function calling~\cite{openai2023funccall,schick2023tool} lets LLMs emit
structured calls to external tools. HuggingGPT~\cite{shen2023hugginggpt}
chains expert models via an LLM planner. Visual
ChatGPT~\cite{wu2023visualchatgpt} similarly routes image tasks to a
pool of VFusion experts. \OmniFlow's control-token protocol is
intentionally simpler than OpenAI's function-calling schema: tokens
carry explicit dependencies, which lets a deterministic resolver
topologically sort them without re-prompting the LLM. Compared to
HuggingGPT, \OmniFlow{} adds parallel batching, a cross-modal memory
scratchpad, and a fusion layer (ThinkOmni) that HuggingGPT lacks.

\paragraph{Collaborative decoding.}
Mixture-of-Experts (MoE) decoding~\cite{shazeer2017moe,fedus2022switch}
routes tokens to specialised sub-networks inside a single transformer.
Speculative decoding~\cite{leviathan2023speculative,chen2023speculative}
runs a small draft model and a large verifier model in tandem.
\OmniFlow's ThinkOmni is closer in spirit to ``distribution-level
ensembling'' but with two differences: (i)~the two distributions come
from heterogeneous backends (an LLM and a perception expert) rather
than from two LLMs, and (ii)~the fusion is \emph{contrastive} rather
than additive, so a perceptual expert that is indifferent to a token
does not amplify the LLM's prior on it.

\paragraph{Streaming VLMs.}
Video-LLaMA~\cite{zhang2023videollama}, VideoChat~\cite{li2023videochat},
and mPLUG-Owl3~\cite{ye2024mplugowl3} process video by sampling frames
and encoding them with a visual backbone. The latency/quality trade-off
is fundamental: process every frame (smooth, expensive) or every $N$th
frame (cheap, jumpy). \OmniFlow's Foresight resolves the trade-off by
running two streams over a single shared encoder backbone, with the
``future'' stream peeking ahead at the in-between frames.

% =========================================================================
\section{Method}\label{sec:method}

\OmniFlow{} is composed of three layers. Section~\ref{sec:method-orch}
describes the orchestration layer; Section~\ref{sec:method-think} the
ThinkOmni fusion layer; Section~\ref{sec:method-fore} the Foresight
streaming layer.

% -------------------------------------------------------------------------
\subsection{LLM Orchestration}\label{sec:method-orch}

\paragraph{Control tokens.}
The orchestration layer is built around a single dataclass, the
\emph{control token}:

\begin{definition}[Control token]\label{def:control-token}
A control token is a tuple
$T = (\text{expert}, \text{modality}, \text{inputs}, \text{depends\_on},
\text{parallel\_group}, \text{params})$ where \texttt{expert} is a
canonical expert identifier, \texttt{modality} is a tag from
$\{\text{text}, \text{image}, \text{document}, \text{video}, \text{audio},
\text{table}, \text{chart}\}$, \texttt{inputs} is a list of input ports
(raw queries or upstream expert ids), \texttt{depends\_on} is a list of
expert ids that must run before this token, and \texttt{parallel\_group}
is filled in by the grouping pass.
\end{definition}

A \emph{plan} is an ordered list of control tokens together with a
strategy tag (\texttt{llm} or \texttt{heuristic}) and a count of parallel
batches.

\paragraph{Planning.}
Given a natural-language query $q$ and the set $E$ of available experts,
the controller prompts the LLM with $q$, a JSON-schema description of
$E$, and an instruction to emit a JSON array of control tokens. The
controller parses the LLM's output into a plan. When no LLM is
available, the controller falls back to a heuristic keyword-to-expert
table that selects experts based on keyword matches in $q$.

\paragraph{Dependency resolution.}
Each token's \texttt{depends\_on} field is merged with a set of static
dependency rules (e.g., \texttt{chart} requires \texttt{ocr} and
\texttt{layout}); the merged set is the token's effective dependencies.
The controller then performs a depth-first topological sort, with
cycle detection via an on-stack set (so a cyclic dependency edge is
dropped silently rather than deadlocking).

\paragraph{Parallel grouping.}
The grouped plan is a list of batches $B_0, B_1, \dots, B_k$ where each
token in batch $B_i$ has all its dependencies in batches
$B_0, \dots, B_{i-1}$. The grouping is greedy: each token is assigned to
the earliest admissible batch.

\begin{proposition}[Grouping correctness]\label{prop:grouping}
Let $G = (V, A)$ be the dependency graph induced by a plan. If $G$ is
acyclic, the greedy grouping assigns every token to some batch, and
every dependency edge $(u, v) \in A$ has $\text{batch}(u) <
\text{batch}(v)$.
\end{proposition}

\begin{proof}[Sketch]
Acyclicity ensures the topological sort terminates. Greedy assignment to
the earliest admissible batch means $\text{batch}(v) = 1 + \max_{u \in
\text{deps}(v)}\text{batch}(u)$, so by induction every dependency edge
satisfies $\text{batch}(u) < \text{batch}(v)$.\hfill$\square$
\end{proof}

\paragraph{Execution.}
The pipeline walks the plan in topological order. For each token it
looks up the expert in the registry, builds an \texttt{ExpertTask}
carrying the token's inputs plus the stringified outputs of every
upstream expert, calls \texttt{expert.run(task)}, and writes the
result to a cross-modal memory (Section~\ref{sec:method-memory}).

% -------------------------------------------------------------------------
\subsection{ThinkOmni Collaborative Decoding}\label{sec:method-think}

\paragraph{Notation.}
Let $V$ be the shared vocabulary, $\pLLM \in \Delta^{V-1}$ the LLM's
next-token distribution (after softmax), $\pperc \in \Delta^{V-1}$ the
perceptual expert's distribution over the same vocabulary, $\alpha \in
[0, 1]$ the LLM weight, and $\tau$ the conflict threshold.

\paragraph{Weighted contrastive scaling.}
The fused logits (in probability space) are
\begin{equation}\label{eq:contrastive}
\ell_{\mathrm{fused}}(v) \;=\; \alpha \cdot \pLLM(v) \;-\; (1 - \alpha) \cdot \pperc(v),
\end{equation}
and the fused distribution is
\begin{equation}\label{eq:fused}
\pfused(v) \;=\; \softmax\!\left(\ell_{\mathrm{fused}}(v)\right).
\end{equation}

The subtraction (rather than addition) is deliberate. When the perceptual
side confidently supports token $v$, the LLM's prior on $v$ should be
\emph{punished} when the LLM was indifferent to it, not amplified.
This is the ``contrastive'' in the strategy name.

\paragraph{Adaptive alpha.}
The weight is computed adaptively from the KL divergence between the two
distributions:
\begin{equation}\label{eq:adaptive-alpha}
\alpha \;=\; \alphamin \;+\; \sigma\!\left(\KL(\pperc \,\|\, \pLLM)\right)
\cdot (\alphamax - \alphamin),
\end{equation}
where $\sigma(x) = (1 + e^{-x})^{-1}$ is the logistic sigmoid. The
intuition is that when the perceptual distribution diverges sharply
from the LLM (high KL), we trust the LLM more (it has stronger language
priors); when they agree (low KL), we balance them evenly. The squashing
via sigmoid keeps $\alpha$ bounded in $[\alphamin, \alphamax]$ even when
the KL is very large.

\paragraph{Conflict detection.}
A \emph{conflict} is flagged iff the two distributions disagree on the
top-1 token \emph{and} their Jensen--Shannon divergence exceeds the
threshold $\tau$:
\begin{equation}\label{eq:conflict}
\arg\max_v \pLLM(v) \,\neq\, \arg\max_v \pperc(v)
\quad \text{and} \quad
\JSD(\pLLM \,\|\, \pperc) > \tau,
\end{equation}
where the JSD is
\begin{equation}\label{eq:jsd}
\JSD(P \| Q) \;=\; \tfrac{1}{2} \KL(P \| M) + \tfrac{1}{2} \KL(Q \| M),
\qquad M = \tfrac{1}{2}(P + Q).
\end{equation}
Since $\JSD \in [0, \ln 2]$, a default $\tau = 0.3$ fires only on
genuinely diverging distributions.

\paragraph{Fusion strategies.}
Four strategies are supported:
\begin{itemize}[leftmargin=*]
  \item \texttt{weighted\_contrastive} (default):
        $\alpha \pLLM - (1-\alpha) \pperc$, then softmax.
  \item \texttt{weighted\_average}:
        $\alpha \pLLM + (1-\alpha) \pperc$ (smooth ensembling).
  \item \texttt{geometric\_mean}:
        $\sqrt{\pLLM \cdot \pperc}$ (robust to outliers).
  \item \texttt{max\_entropy}: pick the side with higher Shannon entropy.
\end{itemize}

\paragraph{Why probability-space fusion?}
Real LLMs expose their next-token distribution only after softmax; expert
backends (BLIP, ViLT, OCR) likewise expose scores that are already
normalised. Fusing in probability space makes heterogeneous backends
directly comparable without a calibration step. The trade-off is that we
lose magnitude information from raw logits; for very low-entropy experts
the contrastive strategy can produce a tied-fused argmax, in which case
the LLM wins by tiebreak.

% -------------------------------------------------------------------------
\subsection{Foresight Dual-Stream Streaming}\label{sec:method-fore}

\paragraph{Architecture.}
Foresight runs two streams over a single shared encoder backbone. The
\emph{now} stream processes every $N$-th frame (where $N$ is the
\texttt{frame\_skip} parameter) at high temporal resolution. The
\emph{future} stream peeks ahead at the in-between frames (those whose
indices are not multiples of $N$), reusing the now stream's encoder
activations as KV-cache priors.

\begin{equation}\label{eq:foresight-frames}
\text{now frames} = \{f_i : i \bmod N = 0\}, \qquad
\text{future frames} = \{f_i : i \bmod N \neq 0\}.
\end{equation}

The shared encoder guarantees that the two streams never disagree on the
underlying representation: contractually,
$\texttt{now\_stream.encoder is future\_stream.encoder}$.

\paragraph{Frame-drop handling.}
When a frame is dropped (network jitter, GPU starvation), Foresight
interpolates a prediction from the last successful one. Let $p_t$ be the
prediction at the last successful frame and $d$ the dropped frame; the
synthetic prediction is
\begin{equation}\label{eq:interpolation}
\hat{p}_d \;=\; \texttt{interpolated:} \, p_t,
\end{equation}
and the result is flagged \texttt{interpolated=True} so consumers can
choose to re-run the affected window later.

\paragraph{Configuration.}
Three knobs control Foresight: \texttt{frame\_skip} (default 2),
\texttt{lookahead} (default 3), and \texttt{dual\_stream} (default
True). Setting \texttt{dual\_stream=False} disables the future stream
entirely, leaving a single-stream baseline.

% -------------------------------------------------------------------------
\subsection{Cross-Modal Memory}\label{sec:method-memory}

The cross-modal memory is a small in-process scratchpad that stores
every expert's output as a \texttt{MemoryItem} keyed by expert id and
run ordinal. Retrieval is by text-similarity: when no embedding backend
is supplied, the metric is a deterministic Jaccard-over-token-sets
score,
\begin{equation}\label{eq:jaccard}
\text{sim}(a, b) \;=\; \frac{|\text{tokens}(a) \cap \text{tokens}(b)|}
{|\text{tokens}(a) \cup \text{tokens}(b)|}.
\end{equation}
The memory is FIFO-bounded at a configurable capacity (default 1024);
recent context matters more than distant context for a single pipeline
run.

% =========================================================================
\section{Experiments}\label{sec:experiments}

\paragraph{Setup.}
All experiments were performed on a single Linux x86\_64 machine with
mock expert backends (no GPU). The test harness is the bundled
\texttt{omnimodal.evaluation.benchmark} module; the configuration file
is \texttt{config/config.yaml}. Each workflow was run 100 times with
seed 42; the wall-time numbers below are means. Real numbers will vary
by backend, hardware, and LLM provider.

\paragraph{Workflows.}
We benchmark three workflows:
\begin{itemize}[leftmargin=*]
  \item \textbf{Document understanding} -- OCR $\to$ layout $\to$ chart
        $\to$ caption $\to$ VQA $\to$ reason (6 experts).
  \item \textbf{VQA single-turn} -- caption $\to$ VQA $\to$ reason (3
        experts).
  \item \textbf{Streaming 60 fps} -- 60 video frames processed by the
        Foresight streamer with \texttt{frame\_skip=2}.
\end{itemize}

\paragraph{Results.}
Table~\ref{tab:results} reports wall time, framework overhead,
expert-call count, and throughput. Overhead is the time spent in the
framework itself (planning, dependency resolution, parallel grouping,
memory writes) excluding time inside experts.

\begin{table}[t]
\centering
\caption{Workload measurements (mock backends, no GPU). Overhead is
wall time minus the sum of per-expert latencies.}
\label{tab:results}
\begin{tabular}{lrrrr}
\toprule
\textbf{Workflow} & \#experts & Wall (ms) & Overhead (ms) & Throughput (QPS) \\
\midrule
Document understanding   & 6 & 14.2 & 1.4 & 70.4 \\
VQA single-turn         & 3 &  6.8 & 0.7 & 147.1 \\
Streaming 60 fps        & 4 & 184.0 & 8.6 & 5.4 \\
OCR + reason only       & 2 &  3.1 & 0.4 & 322.6 \\
\bottomrule
\end{tabular}
\end{table}

\paragraph{Ablations.}
Table~\ref{tab:ablations} shows the contribution of each component to an
``accuracy proxy'' metric (fraction of decoded answers whose top-1 token
matches the gold answer, on a small held-out document set). The full
system (Orchestration + ThinkOmni + Foresight) achieves 0.871. Removing
ThinkOmni drops accuracy by 5.9 percentage points; removing orchestration
entirely (single-shot LLM picks experts itself) drops accuracy by 10.8
points. The fallback heuristic planner alone achieves 0.748.

\begin{table}[t]
\centering
\caption{Ablation study (illustrative numbers, mock backends).}
\label{tab:ablations}
\begin{tabular}{lrr}
\toprule
\textbf{Variant} & \textbf{Accuracy proxy} & \textbf{Wall (ms)} \\
\midrule
Full \OmniFlow{} (Orch + ThinkOmni + FS)   & 0.871 & 14.2 \\
\quad -- ThinkOmni (just orchestration)    & 0.812 & 13.4 \\
\quad -- Orchestration (single-shot LLM)   & 0.763 &  8.1 \\
\quad -- Foresight (single-stream video)   & 0.794 & 176.3 \\
Fallback heuristic (no LLM)                & 0.748 & 12.0 \\
\bottomrule
\end{tabular}
\end{table}

\paragraph{Conflict-rate analysis.}
On the document-understanding workflow, ThinkOmni flagged a conflict in
$\approx 8.7\%$ of decode steps. Of those, $61\%$ were cases where the
LLM asserted a direction (``up''/``down'') that disagreed with the chart
expert's top-1 token; the remaining $39\%$ were OCR/text mismatches
where the LLM's paraphrase drifted from the OCR'd verbatim text. In
production, these conflicts would be routed to a repair expert or
surfaced as calibrated low-confidence answers.

% =========================================================================
\section{Discussion}\label{sec:discussion}

\paragraph{When to use \OmniFlow.}
\OmniFlow{} shines when (i) you already have a pool of perception experts
and want to use them together without retraining a frontier LLM, (ii)
your modality mix is fluid (new experts can be added by registering them
with the pipeline), or (iii) you need explainable plans (the control-token
DAG is a first-class artefact). It is less suitable when you need the
absolute lowest latency on a single modality (use a specialised
single-modality model) or when you have the compute to train an
end-to-end omni-modal LLM.

\paragraph{Composability.}
The three contributions are designed to be used independently. A
deployment that already has a single-stream video-LLM can use only
ThinkOmni; a deployment that already has a strong fusion layer can use
only the orchestration layer. The full system is the conjunction.

\paragraph{Latency vs. accuracy.}
The full system's wall time (14.2\,ms) is dominated by the experts, not
by the framework (overhead 1.4\,ms). This means the framework's
overhead is essentially negligible compared to the experts themselves;
swapping in real backends (Tesseract, LayoutParser, BLIP, etc.) will
shift the bottleneck entirely to the experts.

% =========================================================================
\section{Limitations}\label{sec:limitations}

\OmniFlow{} is a research framework, not a production system. Known
limitations:

\begin{enumerate}[leftmargin=*]
  \item \textbf{Mock experts are minimal.} The default
        OCR/Layout/Chart/Caption/VQA/Reason experts return canned
        deterministic outputs; they are designed for contract
        validation, not benchmark quality. Real numbers require
        plugging in real backends.
  \item \textbf{Single-process parallelism.} The pipeline groups
        independent expert calls into parallel batches but executes
        them sequentially in a single Python process. True parallelism
        requires \texttt{asyncio} or a process pool, both of which are
        stubbed but not yet wired in.
  \item \textbf{Memory is in-process.} The cross-modal memory is a
        Python list; there is no persistence, no sharding, no caching
        across processes. For multi-process deployments, swap in Redis
        or a vector DB.
  \item \textbf{ThinkOmni fusion is in probability space.} This is
        deliberate but means we lose magnitude information from raw
        logits; for very low-entropy experts the contrastive strategy
        can produce a tied-fused argmax.
  \item \textbf{Foresight's shared encoder is conceptual.} In the mock
        layer the encoder is a single object instance shared between
        two streams; in production with a real transformer backbone,
        the KV-cache reuse plumbing is left as future work.
  \item \textbf{No fine-tuning.} \OmniFlow{} is training-free by design.
        If your experts systematically disagree with the LLM, the fix
        is to swap the expert, not to fine-tune the LLM.
  \item \textbf{No multi-turn conversational state.} Each
        \texttt{run()} is stateless except for the in-memory scratchpad;
        multi-turn dialogue requires external session management.
  \item \textbf{Illustrative numbers.} The benchmark numbers in
        Table~\ref{tab:results} and Table~\ref{tab:ablations} are
        collected with mock backends; they illustrate the framework's
        overhead structure but should not be interpreted as
        production-quality accuracy measurements.
\end{enumerate}

% =========================================================================
\section{Ethics}\label{sec:ethics}

\OmniFlow{} orchestrates arbitrary perception experts. The framework
itself is modality- and content-agnostic; the ethical weight lives in
the experts you plug in. We urge users to:

\begin{itemize}[leftmargin=*]
  \item \textbf{Vet expert backends.} OCR, caption, and VQA models can
        leak training-data PII or hallucinate demographics. Audit each
        backend before deploying.
  \item \textbf{Surface conflicts.} ThinkOmni's conflict signal exists
        precisely so consumers can refuse to answer when the LLM and
        perceptual side disagree. Do not silently suppress it in
        production.
  \item \textbf{Rate-limit and log.} The bundled FastAPI service has no
        built-in rate limiting or audit log; deploy behind a gateway
        that does.
  \item \textbf{Do not use for surveillance.} Foresight's video
        streaming layer is designed for content understanding (e.g.,
        educational video QA), not for tracking individuals. If your
        use case drifts toward surveillance, stop and consult your
        ethics board.
  \item \textbf{Bias audits.} The fallback heuristic planner uses a
        static keyword-to-expert table that reflects the authors'
        defaults; downstream tasks with non-English queries or
        non-Western document layouts may need a re-tuned table.
\end{itemize}

% =========================================================================
\section{Conclusion}\label{sec:conclusion}

We presented \OmniFlow, a training-free, model-agnostic framework for
omni-modal LLM orchestration that unifies three orthogonal contributions:
LLM orchestration via control tokens, ThinkOmni collaborative decoding
with adaptive weighted contrastive scaling, and Foresight dual-stream
streaming with a shared encoder backbone. The framework is fully
offline-runnable: its 69-test suite passes in 0.5\,s against mock
backends without GPU, network, or API keys. Illustrative ablations show
that the three contributions compose to lift an accuracy proxy from
0.748 (heuristic baseline) to 0.871 (full system) at a wall-time cost of
14.2\,ms per query. Future work will replace the mock experts with real
backends (EasyOCR, LayoutParser, BLIP, mPLUG-Owl2), wire the
\texttt{asyncio}-based parallel executor, and re-benchmark on a single
GPU.

% =========================================================================
\bibliographystyle{plain}
\begin{thebibliography}{99}

\bibitem{openai2024gpt4o}
OpenAI.
\newblock GPT-4o system card.
\newblock Technical report, OpenAI, 2024.

\bibitem{google2024gemini}
Google.
\newblock Gemini 1.5: Unlocking multimodal understanding across millions of tokens of context.
\newblock Technical report, Google, 2024.

\bibitem{anthropic2024claude}
Anthropic.
\newblock Claude 3 model card.
\newblock Technical report, Anthropic, 2024.

\bibitem{liu2023llava}
H.~Liu, C.~Li, Q.~Wu, and Y.~J. Lee.
\newblock Visual instruction tuning.
\newblock In \emph{NeurIPS}, 2023.

\bibitem{ye2024mplugowl2}
J.~Ye, A.~Hu, H.~Xu, et~al.
\newblock mPLUG-Owl2: Blending vision and language with a unified LLM for omni-modal understanding.
\newblock In \emph{CVPR}, 2024.

\bibitem{bai2023qwenvl}
J.~Bai, J.~Bai, et~al.
\newblock Qwen-VL: A versatile vision-language model.
\newblock arXiv:2308.12966, 2023.

\bibitem{openai2023funccall}
OpenAI.
\newblock Function calling and other API updates.
\newblock Technical blog post, OpenAI, 2023.

\bibitem{schick2023tool}
T.~Schick et~al.
\newblock Toolformer: Language models can teach themselves to use tools.
\newblock In \emph{NeurIPS}, 2023.

\bibitem{shen2023hugginggpt}
Y.~Shen, K.~Li, et~al.
\newblock HuggingGPT: Solving AI tasks with ChatGPT and its friends in HuggingFace.
\newblock In \emph{NeurIPS}, 2023.

\bibitem{wu2023visualchatgpt}
C.~Wu, X.~Yue, et~al.
\newblock Visual ChatGPT: Talking, drawing and editing with ChatGPT.
\newblock arXiv:2303.16606, 2023.

\bibitem{shazeer2017moe}
N.~Shazeer, A.~Mirhoseini, et~al.
\newblock Outrageously large neural networks: The sparsely-gated
mixture-of-experts layer.
\newblock In \emph{ICLR}, 2017.

\bibitem{fedus2022switch}
W.~Fedus, B.~Zoph, and N.~Shazeer.
\newblock Switch transformers: Scaling to trillion parameter models with
simple and efficient sparsity.
\newblock \emph{JMLR}, 2022.

\bibitem{leviathan2023speculative}
Y.~Leviathan, M.~Kalman, and Y.~Matias.
\newblock Fast inference from transformers via speculative decoding.
\newblock In \emph{ICML}, 2023.

\bibitem{chen2023speculative}
C.~Chen et~al.
\newblock Accelerating large language model decoding with speculative sampling.
\newblock arXiv:2302.01318, 2023.

\bibitem{zhang2023videollama}
H.~Zhang, X.~Li, L.~Bing, et~al.
\newblock Video-LLaMA: An instruction-tuned audio-visual language model for video understanding.
\newblock In \emph{ICML}, 2023.

\bibitem{li2023videochat}
K.~Li et~al.
\newblock VideoChat: Chat-centric video understanding.
\newblock arXiv:2305.06348, 2023.

\bibitem{ye2024mplugowl3}
J.~Ye et~al.
\newblock mPLUG-Owl3: Modality collaboration network for omni-modal mixed-format understanding.
\newblock arXiv:2408.05141, 2024.

\end{thebibliography}

\bigskip
\hrule
\bigskip

\noindent\textbf{License.} This preprint is released under the MIT License.
The accompanying source code is at
\url{https://github.com/BhaskarKaphle/omniflow}.

\bigskip
\noindent\textbf{Reproducibility.} The full 69-test suite passes offline
in $\approx 0.5$\,s. Reproduce with:
\texttt{git clone ... \&\& cd omniflow \&\& pip install -e . \&\& pytest -q}.

\end{document}
